\chapter{Phương pháp đề xuất}
\label{Chapter3}

\section{Tổng quan phương pháp đề xuất}
\label{sec:final-pipeline}

Nhóm nghiên cứu đề xuất XBone-Net như một pipeline học sâu đa phương thức dành cho phân loại ảnh X-quang xương, trong đó ảnh cung cấp bằng chứng hình thái và bệnh sử lâm sàng cung cấp ngữ cảnh liên quan đến tình trạng của người bệnh trước khi ảnh được diễn giải. Pipeline kế thừa bộ mã hóa thị giác và bộ mã hóa văn bản của BiomedCLIP \cite{BiomedCLIP}, thích nghi miền dữ liệu cho hai bộ mã hóa bằng LoRA, sử dụng đồng thời thông tin toàn cục và chi tiết cục bộ của ảnh, trao đổi thông tin giữa hai phương thức bằng chú ý chéo hai chiều, sau đó đưa biểu diễn dung hợp vào bộ phân loại dựa trên tâm lớp thực nghiệm. Báo cáo kết quả X-quang không được sử dụng làm đầu vào của pipeline chính nhằm hạn chế nguy cơ mô hình dựa vào từ khóa chẩn đoán thay vì học mối liên hệ giữa hình ảnh và bệnh sử.

Quy trình được chia thành quá trình huấn luyện ngoại tuyến và quá trình suy luận trực tuyến. Cách phân chia này làm rõ thành phần nào được cập nhật từ dữ liệu huấn luyện và thành phần nào chỉ thực hiện truyền thuận khi hệ thống tiếp nhận một ca mới. Trong cả hai quá trình, ảnh và văn bản đi qua cùng một quy tắc tiền xử lý và cùng một kiến trúc mã hóa; khác biệt chủ yếu nằm ở việc tính gradient, cập nhật tham số và ước lượng tâm lớp.

\subsection{Huấn luyện hai giai đoạn trong quá trình ngoại tuyến (Offline)}

Trong giai đoạn thứ nhất, ảnh X-quang và bệnh sử lâm sàng tương ứng được đưa qua hai bộ mã hóa để tạo biểu diễn trong một không gian chung. Mục tiêu khớp ngữ nghĩa được sử dụng để tăng mức tương hợp của cặp ảnh-văn bản tương ứng và ghi nhận quan hệ giữa các mẫu cùng lớp, qua đó giảm việc xem mọi cặp không trùng chỉ số là cặp âm. Mini-batch được lấy mẫu có nhận biết lớp để tăng khả năng mỗi lớp đóng góp nhiều hơn một mẫu vào cùng một batch, nhờ đó quan hệ đồng lớp trong mục tiêu mềm có thể được quan sát trực tiếp. Ở giai đoạn này, trọng số nền của BiomedCLIP được giữ cố định; các adapter LoRA và các thành phần tổng hợp đặc trưng ảnh cục bộ được cập nhật. Bộ dung hợp chú ý chéo và bộ phân loại chưa tham gia vào mục tiêu tối ưu của giai đoạn thứ nhất.

Trong giai đoạn thứ hai, các adapter LoRA được giữ ở dạng tách rời và tiếp tục được cập nhật thay vì hợp nhất vào trọng số nền. Biểu diễn ảnh và văn bản được đưa qua chú ý chéo hai chiều để hình thành biểu diễn dung hợp, sau đó các tâm lớp được ước lượng từ toàn bộ biểu diễn của tập huấn luyện. Hàm mất mát phân loại cập nhật adapter LoRA, thành phần biểu diễn ảnh cục bộ, mô-đun dung hợp và độ lệch riêng của từng lớp; các tâm lớp là thống kê thực nghiệm nên không nhận gradient. Việc tách hai giai đoạn giúp quá trình thích nghi không gian ảnh-văn bản được thực hiện trước khi tối ưu trực tiếp ranh giới quyết định của bài toán phân loại.

Sau mỗi epoch của giai đoạn thứ hai, tâm lớp được ước lượng lại bằng trạng thái hiện tại của bộ mã hóa trước khi đánh giá trên tập validation. Checkpoint được lựa chọn theo độ đo đã xác định trước trên tập validation, sau đó tâm lớp được tính lại một lần cuối từ tập huấn luyện bằng checkpoint được chọn. Tập kiểm thử không tham gia vào quá trình cập nhật tham số, lựa chọn checkpoint hoặc ước lượng tâm lớp.

\subsection{Suy luận trong quá trình trực tuyến (Online)}

Trong quá trình trực tuyến, một ảnh X-quang và bệnh sử lâm sàng của ca cần phân tích được xử lý bằng đúng pipeline đã sử dụng trong giai đoạn thứ hai. Các tham số của bộ mã hóa, adapter LoRA, mô-đun dung hợp, tâm lớp và độ lệch lớp đều được cố định; hệ thống chỉ thực hiện truyền thuận để tạo biểu diễn dung hợp, tính điểm của từng lớp và trả về lớp có điểm lớn nhất. Quá trình này không yêu cầu cập nhật gradient, không ước lượng lại tâm lớp từ mẫu mới và không sử dụng thông tin từ tập kiểm thử để điều chỉnh quyết định.

Biểu diễn dung hợp và đầu ra phân loại có thể được chuyển sang mô-đun OOD hậu xử lý đã hiệu chỉnh bằng dữ liệu huấn luyện và validation. Mô-đun OOD không làm thay đổi dự đoán phân loại bên trong XBone-Net mà bổ sung một tín hiệu về mức độ phù hợp của mẫu với phân phối đã quan sát. Vì vậy, kết quả phân loại và cảnh báo OOD được xem là hai đầu ra có mục đích khác nhau: đầu ra thứ nhất lựa chọn lớp gần nhất trong tập lớp đã biết, còn đầu ra thứ hai xem xét liệu mẫu có đủ tương đồng với dữ liệu trong phân phối hay không.

\section{Các thành phần của pipeline}

\subsection{Đầu vào và tiền xử lý (Input and Preprocessing)}

Mỗi mẫu huấn luyện gồm một ảnh X-quang, một bệnh sử lâm sàng và một nhãn lớp duy nhất. Tập dữ liệu huấn luyện được ký hiệu như sau.

\begin{equation}\mathcal{D}_{\mathrm{train}}=\left\{\left(X_i,R_i,y_i\right)\right\}_{i=1}^{N}.\label{eq:ch3_training_set}\end{equation}

Trong Công thức \ref{eq:ch3_training_set}, $N$ là số mẫu huấn luyện, $X_i$ là ảnh X-quang gốc của mẫu thứ $i$, $R_i$ là bệnh sử lâm sàng tương ứng và $y_i\in\{1,\ldots,K\}$ là nhãn của một trong $K$ lớp. Ký hiệu $\mathcal{D}_{\mathrm{train}}$ chỉ tập huấn luyện và không bao gồm tập validation hoặc tập kiểm thử.

Ảnh X-quang thường có độ phân giải và tỷ lệ khung hình khác nhau, trong khi bộ mã hóa thị giác kế thừa nhận đầu vào có kích thước cố định. Pipeline vì vậy tạo một ảnh toàn cục để giữ bố cục giải phẫu và bốn ảnh cục bộ để giữ lại chi tiết tại các vùng giàu thông tin. Ảnh toàn cục được tạo bằng phép biến đổi giữ tỷ lệ và thêm biên khi cần thiết, sau đó được chuẩn hóa theo quy tắc đầu vào của bộ mã hóa thị giác.

\begin{equation}X_i^{g}=\mathcal{T}_{g}\left(X_i\right)\in\mathbb{R}^{3\times H\times W}.\label{eq:ch3_global_view}\end{equation}

Trong Công thức \ref{eq:ch3_global_view}, $\mathcal{T}_{g}$ là phép tiền xử lý ảnh toàn cục, $X_i^{g}$ là ảnh toàn cục của mẫu thứ $i$, ba kênh đầu vào được tạo theo định dạng mà backbone yêu cầu, còn $H=W=224$ là chiều cao và chiều rộng đưa vào bộ mã hóa. Phép biến đổi này không sử dụng nhãn hoặc vị trí tổn thương.

Đối với nhánh cục bộ, ảnh nguồn trước hết được xác định vùng nội dung nhằm hạn chế các biên đen và vùng đệm của thiết bị chụp. Bốn cửa sổ được chọn theo chiến lược sparse-focal: hai cửa sổ hỗ trợ bao phủ dọc theo trục chính của vùng giải phẫu và hai cửa sổ ưu tiên vùng có năng lượng kết cấu cao, đồng thời hạn chế chồng lấp quá mức với các vùng đã chọn. Mỗi cửa sổ được ánh xạ về ảnh nguồn và biến đổi thành một local view có cùng kích thước đầu vào với ảnh toàn cục.

\begin{equation}\mathcal{X}_i^{\ell}=\left\{X_{i,m}^{\ell}=\mathcal{T}_{\ell}\left(X_i;b_{i,m}\right)\right\}_{m=1}^{M}.\label{eq:ch3_local_views}\end{equation}

Trong Công thức \ref{eq:ch3_local_views}, $\mathcal{X}_i^{\ell}$ là tập local view của mẫu thứ $i$, $M=4$ là số vùng cục bộ cố định, $b_{i,m}$ là hộp tọa độ của vùng thứ $m$ trên ảnh nguồn, $\mathcal{T}_{\ell}$ là phép cắt, đổi kích thước và chuẩn hóa vùng ảnh, còn $X_{i,m}^{\ell}\in\mathbb{R}^{3\times224\times224}$ là local view thứ $m$. Các hộp $b_{i,m}$ chỉ được xác định từ tín hiệu ảnh, không sử dụng nhãn, bounding box tổn thương hoặc một mô hình phát hiện được huấn luyện riêng.

Bệnh sử lâm sàng được chuẩn hóa về mã hóa ký tự, khoảng trắng và các ký hiệu không mang nội dung trước khi token hóa. Chuỗi được cắt hoặc đệm đến độ dài tối đa, đồng thời một attention mask được tạo để phân biệt token nội dung với token đệm.

\begin{equation}\left(\boldsymbol{u}_i,\boldsymbol{m}_i\right)=\mathcal{T}_{t}\left(R_i\right).\label{eq:ch3_text_preprocess}\end{equation}

Trong Công thức \ref{eq:ch3_text_preprocess}, $\mathcal{T}_{t}$ là bộ tiền xử lý và token hóa văn bản, $\boldsymbol{u}_i\in\mathbb{N}^{L}$ là chuỗi chỉ số token, $\boldsymbol{m}_i\in\{0,1\}^{L}$ là attention mask và $L$ là độ dài chuỗi sau khi cắt hoặc đệm. Giá trị một trong $\boldsymbol{m}_i$ biểu thị token hợp lệ, còn giá trị không biểu thị vị trí đệm không được tham gia vào phép chú ý.

\subsection{Bộ mã hóa (Encoders)}

Hai bộ mã hóa được kế thừa từ BiomedCLIP để tận dụng không gian biểu diễn ảnh-văn bản đã được tiền huấn luyện trên dữ liệu y sinh. Bộ mã hóa thị giác xử lý ảnh toàn cục và các local view bằng cùng một tập trọng số, trong khi bộ mã hóa văn bản xử lý chuỗi token của bệnh sử. Đầu ra của hai nhánh được chiếu về cùng chiều $D=512$, nhờ đó có thể được căn chỉnh ở giai đoạn thứ nhất và trao đổi thông tin trong mô-đun dung hợp ở giai đoạn thứ hai.

\subsubsection{Bộ mã hóa thị giác (Visual Encoder)}

Bộ mã hóa thị giác sử dụng Vision Transformer của BiomedCLIP. Ảnh toàn cục được chia thành các patch, kết hợp với token đại diện và embedding vị trí, sau đó đi qua các lớp Transformer. Token đại diện ở đầu ra được chiếu về không gian chung để tạo đặc trưng toàn cục.

\begin{equation}\boldsymbol{g}_i^{I}=\mathcal{P}_{I}\left(\mathcal{E}_{I}\left(X_i^{g}\right)_{\mathrm{CLS}}\right)\in\mathbb{R}^{D}.\label{eq:ch3_global_visual_embedding}\end{equation}

Trong Công thức \ref{eq:ch3_global_visual_embedding}, $\mathcal{E}_{I}$ là bộ mã hóa thị giác, chỉ số $\mathrm{CLS}$ chỉ token đại diện của ảnh, $\mathcal{P}_{I}$ là phép chiếu thị giác sang không gian biểu diễn chung, $\boldsymbol{g}_i^{I}$ là embedding ảnh toàn cục và $D=512$ là số chiều đặc trưng.

Mỗi local view đi qua cùng bộ mã hóa thị giác. Đối với từng vùng, pipeline giữ lại token CLS và một lưới $2\times2$ thu được bằng pooling thích nghi trên các patch token. Cách tổng hợp này tạo năm token cho mỗi vùng, gồm một token mang nội dung tổng quát của vùng và bốn token giữ phân bố không gian thô bên trong vùng.

\begin{equation}\boldsymbol{L}_{i,m}^{I}=\mathcal{P}_{I}\left(\operatorname{Concat}\left[\mathcal{E}_{I}\left(X_{i,m}^{\ell}\right)_{\mathrm{CLS}},\operatorname{Pool}_{2\times2}\left(\mathcal{E}_{I}\left(X_{i,m}^{\ell}\right)_{\mathrm{patch}}\right)\right]\right)\in\mathbb{R}^{5\times D}.\label{eq:ch3_tile_tokens}\end{equation}

Trong Công thức \ref{eq:ch3_tile_tokens}, $\boldsymbol{L}_{i,m}^{I}$ là tập token của local view thứ $m$, ký hiệu $\mathrm{patch}$ chỉ chuỗi patch token ở đầu ra của bộ mã hóa, $\operatorname{Pool}_{2\times2}$ là phép pooling thích nghi tạo bốn vùng không gian, $\operatorname{Concat}$ là phép ghép theo chiều token và $\mathcal{P}_{I}$ là phép chiếu được áp dụng theo từng token. Số chiều của mỗi token sau phép chiếu được ký hiệu là $D$.

Token của bốn local view được ghép thành một chuỗi có ngân sách cố định. Cấu trúc cố định này giúp mô-đun dung hợp nhận cùng số token ảnh đối với mọi mẫu, không phụ thuộc kích thước ban đầu của ảnh X-quang.

\begin{equation}\boldsymbol{L}_{i}^{I}=\operatorname{Concat}_{m=1}^{M}\left(\boldsymbol{L}_{i,m}^{I}\right)\in\mathbb{R}^{P\times D}.\label{eq:ch3_all_local_tokens}\end{equation}

Trong Công thức \ref{eq:ch3_all_local_tokens}, $\boldsymbol{L}_{i}^{I}$ là toàn bộ chuỗi token cục bộ, $M=4$ là số local view và $P=5M=20$ là tổng số token cục bộ. Phép ghép giữ nguyên thứ tự vùng và thứ tự token bên trong từng vùng.

Do các vùng cục bộ được lấy từ những vị trí khác nhau trên ảnh nguồn, tọa độ chuẩn hóa của mỗi vùng được ánh xạ thành embedding không gian và cộng vào token tương ứng. Một hệ số có thể học kiểm soát cường độ của tín hiệu tọa độ để thông tin vị trí bổ sung cho nội dung thị giác mà không thay thế nội dung đó.

\begin{equation}\widetilde{\boldsymbol{L}}_{i}^{I}=\boldsymbol{L}_{i}^{I}+\gamma\,\mathcal{P}_{s}\left(\boldsymbol{B}_{i}\right).\label{eq:ch3_spatial_tokens}\end{equation}

Trong Công thức \ref{eq:ch3_spatial_tokens}, $\boldsymbol{B}_{i}$ chứa tọa độ đã chuẩn hóa của bốn hộp cục bộ, $\mathcal{P}_{s}$ là phép chiếu tọa độ sang không gian $D$ chiều, $\gamma$ là hệ số vô hướng có thể học và $\widetilde{\boldsymbol{L}}_{i}^{I}$ là chuỗi token cục bộ đã bổ sung vị trí. Embedding của một hộp được lặp cho năm token xuất phát từ hộp đó.

Trong giai đoạn thứ nhất, chuỗi token cục bộ được tổng hợp bằng attention pooling có điều kiện theo embedding toàn cục. Cơ chế này cho phép mức đóng góp của các vùng thay đổi theo từng ảnh thay vì sử dụng trung bình không trọng số.

\begin{equation}e_{i,p}=\boldsymbol{w}_{a}^{\mathsf{T}}\tanh\left(\boldsymbol{W}_{g}\boldsymbol{g}_{i}^{I}+\boldsymbol{W}_{\ell}\widetilde{\boldsymbol{L}}_{i,p}^{I}\right).\label{eq:ch3_local_attention_score}\end{equation}

Trong Công thức \ref{eq:ch3_local_attention_score}, $e_{i,p}$ là điểm chú ý chưa chuẩn hóa của token cục bộ thứ $p$, $\widetilde{\boldsymbol{L}}_{i,p}^{I}$ là token cục bộ tương ứng, $\boldsymbol{W}_{g}$ và $\boldsymbol{W}_{\ell}$ là các ma trận chiếu có thể học, $\boldsymbol{w}_{a}$ là vector tham số của attention pooling và $\tanh$ là hàm kích hoạt.

\begin{equation}a_{i,p}=\frac{\exp\left(e_{i,p}\right)}{\sum_{q=1}^{P}\exp\left(e_{i,q}\right)}.\label{eq:ch3_local_attention_weight}\end{equation}

Trong Công thức \ref{eq:ch3_local_attention_weight}, $a_{i,p}$ là trọng số đã chuẩn hóa của token thứ $p$, $P=20$ là tổng số token cục bộ và chỉ số $q$ duyệt qua toàn bộ token cục bộ của cùng một ảnh. Các trọng số không âm và có tổng bằng một.

\begin{equation}\boldsymbol{z}_{i}^{I}=\operatorname{Norm}_{2}\left(\boldsymbol{g}_{i}^{I}+\lambda_{\ell}\sum_{p=1}^{P}a_{i,p}\widetilde{\boldsymbol{L}}_{i,p}^{I}\right).\label{eq:ch3_phase1_visual_embedding}\end{equation}

Trong Công thức \ref{eq:ch3_phase1_visual_embedding}, $\boldsymbol{z}_{i}^{I}$ là embedding thị giác dùng cho mục tiêu căn chỉnh ở giai đoạn thứ nhất, $\operatorname{Norm}_{2}$ là phép chuẩn hóa theo chuẩn Euclid, $\lambda_{\ell}=0{,}25$ là hệ số điều tiết đóng góp của nhánh cục bộ và tổng có trọng số biểu diễn thông tin cục bộ liên quan đến ngữ cảnh toàn cục.

\subsubsection{Bộ mã hóa văn bản (Text Encoder)}

Bộ mã hóa văn bản kế thừa PubMedBERT của BiomedCLIP và nhận chuỗi token cùng attention mask của bệnh sử. Chuỗi đặc trưng đầu ra giữ lại thông tin ở mức token để phục vụ chú ý chéo, trong khi token đại diện được chiếu sang không gian chung để phục vụ căn chỉnh và tạo query toàn cục.

\begin{equation}\boldsymbol{H}_{i}^{T}=\mathcal{P}_{T}\left(\mathcal{E}_{T}\left(\boldsymbol{u}_{i},\boldsymbol{m}_{i}\right)\right)\in\mathbb{R}^{L\times D}.\label{eq:ch3_text_token_embeddings}\end{equation}

Trong Công thức \ref{eq:ch3_text_token_embeddings}, $\mathcal{E}_{T}$ là bộ mã hóa văn bản, $\mathcal{P}_{T}$ là phép chiếu được áp dụng theo từng token, $\boldsymbol{u}_{i}$ là chuỗi chỉ số token, $\boldsymbol{m}_{i}$ là attention mask, $L$ là chiều dài chuỗi và $\boldsymbol{H}_{i}^{T}$ là ma trận biểu diễn theo token sau phép chiếu về $D$ chiều. Các vị trí đệm được loại khỏi phép chú ý bằng attention mask.

\begin{equation}\boldsymbol{z}_{i}^{T}=\operatorname{Norm}_{2}\left(\boldsymbol{H}_{i,0}^{T}\right)\in\mathbb{R}^{D}.\label{eq:ch3_global_text_embedding}\end{equation}

Trong Công thức \ref{eq:ch3_global_text_embedding}, $\boldsymbol{H}_{i,0}^{T}$ là token đại diện ở vị trí đầu chuỗi sau phép chiếu, $\operatorname{Norm}_{2}$ là phép chuẩn hóa $L_2$ và $\boldsymbol{z}_{i}^{T}$ là embedding văn bản toàn cục.

Nhóm nghiên cứu sử dụng LoRA \cite{Hu2022LoRA} để thích nghi các phép chiếu được lựa chọn trong cả hai bộ mã hóa. Trọng số tiền huấn luyện không bị thay thế; phần cập nhật được biểu diễn bằng tích của hai ma trận hạng thấp.

\begin{equation}\boldsymbol{W}'=\boldsymbol{W}+\frac{\alpha_{L}}{r}\boldsymbol{B}\boldsymbol{A}.\label{eq:ch3_lora}\end{equation}

Trong Công thức \ref{eq:ch3_lora}, $\boldsymbol{W}\in\mathbb{R}^{d\times k}$ là ma trận trọng số tiền huấn luyện được giữ cố định, $\boldsymbol{W}'$ là trọng số hiệu dụng khi truyền thuận, $\boldsymbol{A}\in\mathbb{R}^{r\times k}$ và $\boldsymbol{B}\in\mathbb{R}^{d\times r}$ là hai ma trận được huấn luyện, $r$ là hạng LoRA và $\alpha_{L}$ là hệ số tỷ lệ. Điều kiện $r\ll\min(d,k)$ làm cho số tham số cần cập nhật nhỏ hơn đáng kể so với tinh chỉnh toàn bộ ma trận $\boldsymbol{W}$.

Trong giai đoạn thứ nhất, độ tương đồng giữa embedding ảnh và embedding văn bản được tính bằng cosine. Do hai embedding đã được chuẩn hóa, tích vô hướng tương ứng với độ tương đồng cosine; công thức tổng quát vẫn được trình bày để làm rõ đại lượng được tối ưu.

\begin{equation}s_{ij}=\frac{\left(\boldsymbol{z}_{i}^{I}\right)^{\mathsf{T}}\boldsymbol{z}_{j}^{T}}{\left\|\boldsymbol{z}_{i}^{I}\right\|_{2}\left\|\boldsymbol{z}_{j}^{T}\right\|_{2}}.\label{eq:ch3_cross_modal_similarity}\end{equation}

Trong Công thức \ref{eq:ch3_cross_modal_similarity}, $s_{ij}$ là mức tương đồng giữa ảnh thứ $i$ và bệnh sử thứ $j$, $\boldsymbol{z}_{i}^{I}$ là embedding thị giác, $\boldsymbol{z}_{j}^{T}$ là embedding văn bản, ký hiệu $\mathsf{T}$ biểu thị chuyển vị và $\|\cdot\|_{2}$ là chuẩn Euclid.

Mục tiêu mềm phân bổ trọng số lớn cho cặp tương ứng và một phần trọng số cho các cặp khác thuộc cùng lớp. Phân phối mục tiêu theo hàng được xây dựng như sau.

\begin{equation}q_{ij}=\frac{\rho\,\mathbb{1}\left[i=j\right]+\left(1-\rho\right)\mathbb{1}\left[y_i=y_j\right]}{\sum_{n=1}^{B}\left(\rho\,\mathbb{1}\left[i=n\right]+\left(1-\rho\right)\mathbb{1}\left[y_i=y_n\right]\right)}.\label{eq:ch3_semantic_target}\end{equation}

Trong Công thức \ref{eq:ch3_semantic_target}, $q_{ij}$ là xác suất mục tiêu của cặp ảnh $i$ và văn bản $j$, $B$ là kích thước mini-batch, $\rho\in[0,1]$ là hệ số ưu tiên cặp tương ứng, $\mathbb{1}[\cdot]$ là hàm chỉ thị và $y_i,y_j$ là nhãn lớp. Mẫu số chuẩn hóa tổng trọng số mục tiêu của mỗi ảnh bằng một; khi trong mini-batch có nhiều mẫu cùng lớp, các mẫu này nhận trọng số mềm thay vì bị xem hoàn toàn là cặp âm.

Xác suất dự đoán theo hướng ảnh sang văn bản được tính từ toàn bộ văn bản trong cùng mini-batch.

\begin{equation}p_{ij}^{I\rightarrow T}=\frac{\exp\left(s_{ij}/\tau\right)}{\sum_{n=1}^{B}\exp\left(s_{in}/\tau\right)}.\label{eq:ch3_image_to_text_probability}\end{equation}

Trong Công thức \ref{eq:ch3_image_to_text_probability}, $p_{ij}^{I\rightarrow T}$ là xác suất văn bản $j$ tương ứng với ảnh $i$, $s_{ij}$ là độ tương đồng ảnh-văn bản, $\tau>0$ là nhiệt độ và $B$ là kích thước mini-batch. Nhiệt độ nhỏ tạo phân phối sắc hơn, còn nhiệt độ lớn tạo phân phối phẳng hơn.

Xác suất theo hướng văn bản sang ảnh được tính đối xứng trên toàn bộ ảnh trong mini-batch.

\begin{equation}p_{ji}^{T\rightarrow I}=\frac{\exp\left(s_{ij}/\tau\right)}{\sum_{n=1}^{B}\exp\left(s_{nj}/\tau\right)}.\label{eq:ch3_text_to_image_probability}\end{equation}

Trong Công thức \ref{eq:ch3_text_to_image_probability}, $p_{ji}^{T\rightarrow I}$ là xác suất ảnh $i$ tương ứng với văn bản $j$, mẫu số duyệt qua các ảnh có chỉ số $n$ trong mini-batch, còn $s_{ij}$, $\tau$ và $B$ có ý nghĩa như trong Công thức \ref{eq:ch3_image_to_text_probability}.

Hàm mất mát khớp ngữ nghĩa lấy trung bình hai hướng để tránh ưu tiên riêng nhánh ảnh hoặc nhánh văn bản.

\begin{equation}\mathcal{L}_{\mathrm{SM}}=-\frac{1}{2B}\sum_{i=1}^{B}\sum_{j=1}^{B}q_{ij}\left(\log p_{ij}^{I\rightarrow T}+\log p_{ji}^{T\rightarrow I}\right).\label{eq:ch3_semantic_matching_loss}\end{equation}

Trong Công thức \ref{eq:ch3_semantic_matching_loss}, $\mathcal{L}_{\mathrm{SM}}$ là hàm mất mát của giai đoạn thứ nhất, $q_{ij}$ là mục tiêu mềm, $p_{ij}^{I\rightarrow T}$ và $p_{ji}^{T\rightarrow I}$ là hai phân phối dự đoán, còn $B$ là kích thước mini-batch. Gradient của hàm mất mát này cập nhật adapter LoRA và thành phần tổng hợp token ảnh cục bộ nhưng không cập nhật mô-đun dung hợp hoặc bộ phân loại.

\subsection{Bộ dung hợp (Cross-Attention)}

Bộ dung hợp chỉ được tối ưu và sử dụng từ giai đoạn thứ hai. Đầu vào thị giác gồm token toàn cục cùng 20 token cục bộ đã bổ sung tọa độ, trong khi đầu vào văn bản gồm embedding toàn cục và chuỗi token bệnh sử. Hai hướng chú ý có vai trò bổ sung: hướng ảnh sang văn bản dùng ngữ cảnh lâm sàng để làm giàu token ảnh toàn cục, còn hướng văn bản sang ảnh dùng bằng chứng thị giác toàn cục và cục bộ để làm giàu token văn bản toàn cục.

\begin{equation}\boldsymbol{V}_{i}=\operatorname{Concat}\left[\boldsymbol{g}_{i}^{I},\widetilde{\boldsymbol{L}}_{i}^{I}\right]\in\mathbb{R}^{(P+1)\times D}.\label{eq:ch3_visual_sequence}\end{equation}

Trong Công thức \ref{eq:ch3_visual_sequence}, $\boldsymbol{V}_{i}$ là chuỗi token thị giác dùng cho dung hợp, $\boldsymbol{g}_{i}^{I}$ là token ảnh toàn cục, $\widetilde{\boldsymbol{L}}_{i}^{I}$ là $P=20$ token cục bộ đã bổ sung tọa độ và $D=512$ là chiều đặc trưng. Vì vậy, chuỗi thị giác có tổng cộng $P+1=21$ token.

Ở hướng ảnh sang văn bản, token ảnh toàn cục đóng vai trò query, còn chuỗi token văn bản đóng vai trò key và value. Attention mask của văn bản được áp dụng trước Softmax để token đệm không nhận trọng số.

\begin{equation}\boldsymbol{A}_{i}^{I\leftarrow T}=\operatorname{Softmax}\left(\frac{\left(\boldsymbol{g}_{i}^{I}\boldsymbol{W}_{Q}^{I}\right)\left(\boldsymbol{H}_{i}^{T}\boldsymbol{W}_{K}^{T}\right)^{\mathsf{T}}}{\sqrt{d_{h}}}+\boldsymbol{M}_{i}^{T}\right).\label{eq:ch3_image_queries_text_attention}\end{equation}

Trong Công thức \ref{eq:ch3_image_queries_text_attention}, $\boldsymbol{A}_{i}^{I\leftarrow T}$ là phân bố chú ý từ token ảnh đến các token văn bản, $\boldsymbol{W}_{Q}^{I}$ và $\boldsymbol{W}_{K}^{T}$ là ma trận chiếu query và key, $d_h$ là chiều của mỗi attention head, $\boldsymbol{M}_{i}^{T}$ là mặt nạ cộng có giá trị không tại token hợp lệ và giá trị âm lớn tại token đệm, còn $\operatorname{Softmax}$ chuẩn hóa trọng số theo chiều token văn bản.

\begin{equation}\boldsymbol{h}_{i}^{I}=\operatorname{LN}\left(\boldsymbol{g}_{i}^{I}+\boldsymbol{A}_{i}^{I\leftarrow T}\left(\boldsymbol{H}_{i}^{T}\boldsymbol{W}_{V}^{T}\right)\right).\label{eq:ch3_text_enhanced_image}\end{equation}

Trong Công thức \ref{eq:ch3_text_enhanced_image}, $\boldsymbol{h}_{i}^{I}$ là biểu diễn ảnh đã được bổ sung ngữ cảnh văn bản, $\boldsymbol{W}_{V}^{T}$ là ma trận chiếu value của nhánh văn bản, $\operatorname{LN}$ là Layer Normalization và phép cộng với $\boldsymbol{g}_{i}^{I}$ tạo kết nối residual giúp giữ thông tin ảnh ban đầu.

Ở hướng văn bản sang ảnh, embedding văn bản toàn cục đóng vai trò query, còn chuỗi 21 token thị giác đóng vai trò key và value. Hướng này cho phép một biểu diễn văn bản tổng quát lựa chọn thông tin từ bố cục toàn cục và các vùng ảnh cục bộ.

\begin{equation}\boldsymbol{A}_{i}^{T\leftarrow I}=\operatorname{Softmax}\left(\frac{\left(\boldsymbol{z}_{i}^{T}\boldsymbol{W}_{Q}^{T}\right)\left(\boldsymbol{V}_{i}\boldsymbol{W}_{K}^{I}\right)^{\mathsf{T}}}{\sqrt{d_{h}}}\right).\label{eq:ch3_text_queries_image_attention}\end{equation}

Trong Công thức \ref{eq:ch3_text_queries_image_attention}, $\boldsymbol{A}_{i}^{T\leftarrow I}$ là phân bố chú ý từ embedding văn bản đến chuỗi token thị giác, $\boldsymbol{W}_{Q}^{T}$ và $\boldsymbol{W}_{K}^{I}$ là các ma trận chiếu query và key, $\boldsymbol{V}_{i}$ là chuỗi thị giác và $d_h$ là chiều của mỗi attention head.

\begin{equation}\boldsymbol{h}_{i}^{T}=\operatorname{LN}\left(\boldsymbol{z}_{i}^{T}+\boldsymbol{A}_{i}^{T\leftarrow I}\left(\boldsymbol{V}_{i}\boldsymbol{W}_{V}^{I}\right)\right).\label{eq:ch3_image_enhanced_text}\end{equation}

Trong Công thức \ref{eq:ch3_image_enhanced_text}, $\boldsymbol{h}_{i}^{T}$ là biểu diễn văn bản đã được bổ sung thông tin thị giác, $\boldsymbol{W}_{V}^{I}$ là ma trận chiếu value của nhánh ảnh, $\operatorname{LN}$ là Layer Normalization và kết nối residual giữ lại nội dung của embedding văn bản toàn cục.

Hai biểu diễn tăng cường được ghép và đưa qua MLP để tạo embedding dung hợp dùng chung cho bộ phân loại. MLP giảm vector ghép từ $2D=1024$ chiều về $D=512$ chiều, sau đó tiếp tục chiếu trong không gian $D$ chiều và chuẩn hóa đầu ra. Các embedding toàn cục thô không được ghép lại thêm một lần nữa sau bước này.

\begin{equation}\boldsymbol{z}_{i}^{F}=\operatorname{LN}\left(\boldsymbol{W}_{2}\,\phi\left(\boldsymbol{W}_{1}\operatorname{Concat}\left[\boldsymbol{h}_{i}^{I},\boldsymbol{h}_{i}^{T}\right]+\boldsymbol{b}_{1}\right)+\boldsymbol{b}_{2}\right)\in\mathbb{R}^{D}.\label{eq:ch3_fused_embedding}\end{equation}

Trong Công thức \ref{eq:ch3_fused_embedding}, $\boldsymbol{z}_{i}^{F}$ là embedding dung hợp, $\boldsymbol{W}_{1}$ và $\boldsymbol{W}_{2}$ là hai ma trận trọng số của MLP, $\boldsymbol{b}_{1}$ và $\boldsymbol{b}_{2}$ là các vector độ lệch, $\phi$ là hàm kích hoạt phi tuyến, $\operatorname{Concat}$ là phép ghép theo chiều đặc trưng và $D=512$ là chiều đầu ra.

\subsection{Bộ phân loại dựa trên tâm lớp thực nghiệm (Classifier)}

Bộ phân loại biểu diễn mỗi lớp bằng trung bình embedding dung hợp của các mẫu huấn luyện thuộc lớp đó. Tâm lớp không phải là tham số tự do và không được optimizer cập nhật; khi bộ mã hóa hoặc bộ dung hợp thay đổi, các tâm được ước lượng lại từ tập huấn luyện. Thiết kế này giữ cho đại diện lớp gắn trực tiếp với hình học của không gian biểu diễn hiện tại.

\begin{equation}\boldsymbol{c}_{k}=\frac{1}{N_{k}}\sum_{i:y_i=k}\boldsymbol{z}_{i}^{F}.\label{eq:ch3_empirical_centroid}\end{equation}

Trong Công thức \ref{eq:ch3_empirical_centroid}, $\boldsymbol{c}_{k}\in\mathbb{R}^{D}$ là tâm thực nghiệm của lớp $k$, $N_k$ là số mẫu huấn luyện thuộc lớp đó, $\boldsymbol{z}_{i}^{F}$ là embedding dung hợp của mẫu thứ $i$ và điều kiện $y_i=k$ giới hạn phép tổng trên các mẫu có nhãn $k$. Mỗi lớp trong cấu hình phải có ít nhất một mẫu huấn luyện để tâm lớp được xác định.

Độ tương đồng giữa một mẫu và từng tâm lớp được tính bằng cosine, nhờ đó quyết định phụ thuộc chủ yếu vào hướng của embedding thay vì độ lớn tuyệt đối của vector.

\begin{equation}r_{ik}=\frac{\left(\boldsymbol{z}_{i}^{F}\right)^{\mathsf{T}}\boldsymbol{c}_{k}}{\left\|\boldsymbol{z}_{i}^{F}\right\|_{2}\left\|\boldsymbol{c}_{k}\right\|_{2}}.\label{eq:ch3_centroid_similarity}\end{equation}

Trong Công thức \ref{eq:ch3_centroid_similarity}, $r_{ik}$ là độ tương đồng cosine giữa mẫu $i$ và lớp $k$, $\boldsymbol{z}_{i}^{F}$ là embedding dung hợp, $\boldsymbol{c}_{k}$ là tâm lớp, ký hiệu $\mathsf{T}$ biểu thị chuyển vị và $\|\cdot\|_2$ là chuẩn Euclid.

Điểm phân loại kết hợp độ tương đồng với một hệ số tỷ lệ cố định và một độ lệch riêng cho từng lớp. Độ lệch lớp được khởi tạo bằng không và được học trong giai đoạn thứ hai, qua đó cho phép điều chỉnh ranh giới quyết định mà không biến tâm thực nghiệm thành prototype tham số.

\begin{equation}\ell_{ik}=\alpha_{c}\,r_{ik}+b_{k}.\label{eq:ch3_classifier_logit}\end{equation}

Trong Công thức \ref{eq:ch3_classifier_logit}, $\ell_{ik}$ là logit của mẫu $i$ đối với lớp $k$, $r_{ik}$ là độ tương đồng với tâm lớp, $\alpha_c=15$ là hệ số tỷ lệ cố định và $b_k$ là class-specific bias có thể học của lớp $k$.

Logit của tất cả các lớp được chuẩn hóa bằng Softmax để tạo phân phối xác suất dự đoán.

\begin{equation}p_{ik}=\frac{\exp\left(\ell_{ik}\right)}{\sum_{j=1}^{K}\exp\left(\ell_{ij}\right)}.\label{eq:ch3_class_probability}\end{equation}

Trong Công thức \ref{eq:ch3_class_probability}, $p_{ik}$ là xác suất dự đoán lớp $k$ cho mẫu $i$, $\ell_{ik}$ là logit tương ứng, $K$ là tổng số lớp và chỉ số $j$ duyệt qua toàn bộ lớp trong bài toán.

Để hạn chế việc các lớp có nhiều mẫu chi phối hàm mục tiêu, trọng số lớp được xây dựng từ số mẫu hiệu dụng \cite{Cui2019ClassBalanced}. Đại lượng số mẫu hiệu dụng của lớp $k$ được xác định như sau.

\begin{equation}E_{k}=\frac{1-\beta^{N_k}}{1-\beta}.\label{eq:ch3_effective_number}\end{equation}

Trong Công thức \ref{eq:ch3_effective_number}, $E_k$ là số mẫu hiệu dụng của lớp $k$, $N_k$ là số mẫu thực tế của lớp và $\beta\in[0,1)$ là hệ số kiểm soát mức bão hòa. Khi số mẫu tăng, đóng góp bổ sung của các mẫu mới giảm dần thay vì tăng tuyến tính không giới hạn.

Nghịch đảo số mẫu hiệu dụng được chuẩn hóa để trọng số trung bình của các lớp bằng một, đồng thời có thể được chặn trên nhằm hạn chế gradient quá lớn ở lớp rất hiếm.

\begin{equation}w_k=\frac{K\,E_k^{-1}}{\sum_{j=1}^{K}E_j^{-1}}.\label{eq:ch3_class_weight}\end{equation}

Trong Công thức \ref{eq:ch3_class_weight}, $w_k$ là trọng số của lớp $k$, $E_k^{-1}$ là nghịch đảo số mẫu hiệu dụng và $K$ là tổng số lớp. Việc chuẩn hóa không làm các lớp có tần suất bằng nhau, mà chỉ thay đổi mức đóng góp tương đối của chúng vào hàm mất mát.

Label smoothing được áp dụng để tránh mục tiêu one-hot tuyệt đối và giảm mức quá tự tin của đầu phân loại. Phân phối mục tiêu sau làm trơn được xác định như sau.

\begin{equation}\widetilde{y}_{ik}=\left(1-\varepsilon\right)\mathbb{1}\left[y_i=k\right]+\frac{\varepsilon}{K}.\label{eq:ch3_label_smoothing}\end{equation}

Trong Công thức \ref{eq:ch3_label_smoothing}, $\widetilde{y}_{ik}$ là xác suất mục tiêu của lớp $k$ đối với mẫu $i$, $\varepsilon=0{,}05$ là hệ số label smoothing, $\mathbb{1}[y_i=k]$ là hàm chỉ thị của nhãn đúng và $K$ là tổng số lớp.

Hàm mất mát của giai đoạn thứ hai là Cross-Entropy có trọng số lớp trên phân phối mục tiêu đã làm trơn. Pipeline không bổ sung Prototypical Loss dạng pull-push vì các tâm lớp được xác định bằng thống kê của tập huấn luyện thay vì được học như tham số tự do.

\begin{equation}\mathcal{L}_{\mathrm{P2}}=-\frac{1}{B}\sum_{i=1}^{B}\sum_{k=1}^{K}w_{k}\widetilde{y}_{ik}\log p_{ik}.\label{eq:ch3_phase2_loss}\end{equation}

Trong Công thức \ref{eq:ch3_phase2_loss}, $\mathcal{L}_{\mathrm{P2}}$ là hàm mất mát của giai đoạn thứ hai, $B$ là kích thước mini-batch, $w_k$ là trọng số của lớp $k$, $\widetilde{y}_{ik}$ là mục tiêu đã làm trơn và $p_{ik}$ là xác suất dự đoán. Gradient của hàm mất mát cập nhật các tham số có thể học của pipeline nhưng không cập nhật trực tiếp các tâm $\boldsymbol{c}_k$.

Trong quá trình trực tuyến, lớp dự đoán là lớp có logit lớn nhất. Quy tắc này tương đương lựa chọn tâm lớp có điểm cosine đã hiệu chỉnh lớn nhất.

\begin{equation}\widehat{y}_{i}=\underset{k\in\{1,\ldots,K\}}{\operatorname{arg\,max}}\;\ell_{ik}.\label{eq:ch3_prediction}\end{equation}

Trong Công thức \ref{eq:ch3_prediction}, $\widehat{y}_{i}$ là nhãn dự đoán của mẫu $i$, $k$ là chỉ số lớp, $K$ là tổng số lớp và $\ell_{ik}$ là logit đã kết hợp độ tương đồng tâm lớp với class-specific bias. Quy tắc dự đoán sử dụng các tâm và độ lệch lớp đã được cố định sau quá trình huấn luyện ngoại tuyến.

\section{Tóm tắt chương}

Chương này trình bày XBone-Net theo một luồng thống nhất từ đầu vào đến quyết định phân loại. Trong quá trình ngoại tuyến, giai đoạn thứ nhất thích nghi hai bộ mã hóa và nhánh ảnh cục bộ bằng mục tiêu khớp ngữ nghĩa, còn giai đoạn thứ hai tiếp tục thích nghi biểu diễn đồng thời học chú ý chéo, MLP dung hợp và độ lệch lớp dưới mục tiêu Cross-Entropy có trọng số. Trong quá trình trực tuyến, cùng pipeline tiền xử lý, mã hóa và dung hợp được sử dụng với toàn bộ tham số đã cố định; nhãn được xác định bằng khoảng cách tới các tâm lớp thực nghiệm. Cấu trúc này cho phép tách rõ vai trò của bảo toàn chi tiết ảnh, mã hóa bệnh sử, tương tác liên phương thức và ra quyết định dựa trên hình học của tập huấn luyện.



% của chap3 new

\chapter{Phương pháp đề xuất}
\label{Chapter3}

\section{Phân loại zero-shot với BioMedCLIP và định hướng cải tiến}
\label{sec:ch3_biomedclip_baseline}

Với việc được tiền huấn luyện trên PMC-15M, một tập dữ liệu y sinh quy mô lớn có bao gồm ảnh X-quang, BioMedCLIP \cite{BiomedCLIP} được lựa chọn làm mô hình nền tảng để phát triển hệ thống phân loại. Trước khi xây dựng phương pháp đề xuất, nhóm khảo sát khả năng chuyển trực tiếp tri thức của mô hình sang bài toán phân loại ảnh X-quang xương trong thiết lập \textit{suy luận không mẫu}. Khảo sát này giúp làm rõ cơ chế phân loại ban đầu của BioMedCLIP và xác định những khía cạnh cần thích nghi với dữ liệu đích.

Phần này tập trung mô tả quy trình phân loại và các định hướng cải tiến. Thiết kế thực nghiệm được trình bày tại Mục \ref{sec:experiment_content}, còn kết quả định lượng và phần bàn luận tương ứng được trình bày tại Mục \ref{subsec:foundation_model_evaluation} của Chương 4 để tránh lặp lại nội dung thực nghiệm trong chương phương pháp.

\subsection{Quá trình phân loại của BioMedCLIP}
\label{subsec:ch3_biomedclip_classification}

BioMedCLIP trong thiết lập \textit{không mẫu} thực hiện phân loại dựa trên cơ chế đối sánh ảnh-văn bản thay vì sử dụng một đầu phân loại được huấn luyện trên dữ liệu đích. Quy trình gồm năm bước chính:

\begin{enumerate}
    \item \textbf{Xây dựng câu nhắc văn bản:} Với mỗi lớp, tên bệnh lý được đưa vào một câu nhắc mô tả ảnh X-quang, chẳng hạn \texttt{A radiograph showing osteosarcoma.} đối với lớp sarcoma xương. Cùng một mẫu câu được áp dụng nhất quán cho toàn bộ tập lớp.
    \item \textbf{Trích xuất đặc trưng văn bản:} Câu nhắc của từng lớp được đưa qua nhánh mã hóa văn bản sử dụng kiến trúc PubMedBERT để tạo các vectơ đặc trưng lớp trong không gian biểu diễn chung.
    \item \textbf{Trích xuất đặc trưng hình ảnh:} Ảnh X-quang đầu vào được xử lý qua nhánh mã hóa thị giác sử dụng kiến trúc Vision Transformer để tạo vectơ đặc trưng đại diện cho thông tin thị giác.
    \item \textbf{Tính độ tương đồng:} Độ tương đồng côsin được tính giữa vectơ ảnh và vectơ văn bản của từng lớp. Cặp ảnh - văn bản có sự tương hợp cao nhất về mặt ngữ nghĩa y khoa sẽ đạt điểm số lớn nhất.
    \item \textbf{Đưa ra dự đoán:} Các điểm số tương đồng được chuẩn hóa thông qua hàm Softmax để chuyển thành phân phối xác suất. Lớp tương ứng với câu mô tả có xác suất cao nhất sẽ được lựa chọn làm kết quả dự đoán của mô hình.
\end{enumerate}